% Volume X: Normative Conclusions
% Part XVIII. Freedom Under Causal Opacity
% Chapter 109: Disciplined Paranoia
% Spine: proof-spines/ch109-spine.md

\chapter{Disciplined Paranoia}
\label{ch:ch109}

Disciplined paranoia begins from five standing assumptions: every record may be incomplete, every observer situated, every model contestable, every institution fallible, and every concentration of power vulnerable to capture. Those assumptions widen inquiry by refusing to equate the current account with the whole truth.

But disciplined paranoia is not permission for coercion by hunch. It requires specific evidence before specific coercion. Its formula is
\[
\text{maximum permission to question}
\quad+\quad
\text{strict limits on acting from suspicion}.
\]
It is therefore a civic ethic of expansive doubt joined to severe restraint on what doubt alone may authorize.
